% Intended LaTeX compiler: pdflatex
\documentclass[a4paper,12pt]{article}
\usepackage[utf8]{inputenc}
\usepackage[T1]{fontenc}
\usepackage{graphicx}
\usepackage{longtable}
\usepackage{wrapfig}
\usepackage{rotating}
\usepackage[normalem]{ulem}
\usepackage{amsmath}
\usepackage{amssymb}
\usepackage{capt-of}
\usepackage{hyperref}
\usepackage[margin=0.2in]{geometry}
\usepackage{amsmath}
\usepackage{amsfonts}
\usepackage{indentfirst}
\usepackage{pgffor}
\usepackage{amssymb}
\usepackage{cancel}
\usepackage{amsthm}
\usepackage{polynom}
\usepackage{tikz}
\usepackage[tikz]{bclogo}
\usepackage{listings}
\usepackage[most]{tcolorbox}
\usepackage{forest}
\usepackage{adjustbox}
\usepackage{tikz-3dplot}
\usepackage{mathtools}
\usepackage{pgfplots}
\usetikzlibrary{tikzmark,calc,fit,matrix,arrows,automata,positioning}
\usepackage{centernot}
\usepackage{lmodern}
\usepackage{physics}
\usepackage[boxed,linesnumbered,vlined]{algorithm2e}
\usepackage{algpseudocode}
\usepackage{algorithmicx}
\usepackage{svg}
\tikzstyle{every picture}+=[remember picture]
\DeclareMathOperator{\spn}{span}
\DeclareMathOperator{\dom}{domain}
\DeclareMathOperator{\ran}{range}
\DeclareMathOperator{\rng}{range}
\DeclareMathOperator{\img}{Im}
\DeclareMathOperator{\adj}{adj}
\newcommand\dif[1]{\:\textrm{d}#1}
\DeclarePairedDelimiter\ceil{\lceil}{\rceil}
\DeclarePairedDelimiter\floor{\lfloor}{\rfloor}
\newcommand{\ihat}{\hat{\textbf{\i}}}
\newcommand{\jhat}{\hat{\textbf{\j}}}
\newcommand{\khat}{\hat{\textbf{k}}}
\newcommand{\rhat}{\hat{\textbf{r}}}
\newcommand{\thetahat}{\boldsymbol{\hat{\theta}}}
\author{mahmood}
\date{\today}
\title{algorithms homework 3}
\hypersetup{
 pdfauthor={mahmood},
 pdftitle={algorithms homework 3},
 pdfkeywords={},
 pdfsubject={third homework of the algorithms course, on the subject of recursion},
 pdfcreator={Emacs 28.2 (Org mode 9.5.5)}, 
 pdflang={English}}
\begin{document}

\maketitle
\definecolor{Brown}{cmyk}{0,0.81,1,0.60}
\definecolor{OliveGreen}{cmyk}{0.64,0,0.95,0.40}
\definecolor{CadetBlue}{cmyk}{0.62,0.57,0.23,0}
\definecolor{lightlightgray}{gray}{0.9}
\lstset{
  language=C,                             % Code langugage
  basicstyle=\ttfamily,                   % Code font, Examples: \footnotesize, \ttfamily
  keywordstyle=\color{OliveGreen},        % Keywords font ('*' = uppercase)
  commentstyle=\color{gray},              % Comments font
  numbers=left,                           % Line nums position
  numberstyle=\tiny,                      % Line-numbers fonts
  stepnumber=1,                           % Step between two line-numbers
  numbersep=5pt,                          % How far are line-numbers from code
  backgroundcolor=\color{lightlightgray}, % Choose background color
  frame=single,                           % A frame around the code
  tabsize=2,                              % Default tab size
  captionpos=b,                           % Caption-position = bottom
  breaklines=true,                        % Automatic line breaking?
  breakatwhitespace=false,                % Automatic breaks only at whitespace?
  showspaces=false,                       % Dont make spaces visible
  showtabs=false,                         % Dont make tabls visible
  columns=flexible,                       % Column format
  morekeywords={__global__, __device__},  % CUDA specific keywords
  showstringspaces=false,                 % dont show spaces in strings
}
% so that latex doesnt complain about sage not being a language
\lstdefinelanguage{sage}{}

% environment definition for special blocks
\theoremstyle{definition}
\newtheorem{my_example}{Example}
\counterwithin{my_example}{section}
\counterwithin{my_example}{subsection}
\newtheorem{definition}{Definition}
\counterwithin{definition}{section}
\counterwithin{definition}{subsection}
\newtheorem{characteristic}{Characteristic}
\newtheorem{entailment}{Entailment}
%\newtheore*{proof}{Proof}
\newtheorem{lemma}{Lemma}
\newtheorem{question}{Question}
\newtheorem{subquestion}{Subquestion}[question]
%\newtheorem{note}{Note}
\newtheorem{answer}{Answer}
\counterwithin{answer}{question}
\counterwithin{answer}{subquestion}
\newtheorem{step}{Step}[answer]

% the following snippet auto resizes images to fit properly
% Determine if the image is too wide for the page.
% \makeatletter
% \def\ScaleIfNeeded{%
%   \ifdim\Gin@nat@width>\linewidth
%     \linewidth
%   \else
%     \Gin@nat@width
%   \fi
% }
% \makeatother
% % Resize figures that are too wide for the page.
% \let\oldincludegraphics\includegraphics
% \renewcommand\includegraphics[2][]{%
%   \oldincludegraphics[width=\ScaleIfNeeded]{#2}
% }

\newenvironment{note}
  {\begin{bclogo}[logo=\bcinfo, couleurBarre=orange, couleur=lightgray]{ note}}
  {\end{bclogo}}
\newenvironment{attention}
  {\begin{bclogo}[logo=\bcattention, couleurBarre=red, noborder=true, 
                couleur=red!15]{Important!}}
  {\end{bclogo}}

homework on \href{20221105001640-recursive_function.org}{recursion}\\
\begin{note}
please do provide feedback on mistakes you find, if any\\
\end{note}
\begin{question}
write a recursive algorithm that finds the 2 biggest numbers in a given array that halves the array on each recursive call\\
\begin{answer}
\begin{center}
\includesvg{/home/mahmooz/brain/out/HuER61P}
\end{center}

python version of the algorithm:\\
\lstset{language=Python,label= ,caption= ,captionpos=b,numbers=none}
\begin{lstlisting}
def biggest_two(arr):
    if len(arr) == 1:
        return arr[0], -1000
    if len(arr) == 2:
        return arr[0], arr[1]
    left_1, left_2 = biggest_two(arr[0:len(arr)//2])
    right_1, right_2 = biggest_two(arr[len(arr)//2:len(arr)])
    biggest_1 = max({left_1, left_2, right_1, right_2})
    biggest_2 = max({left_1, left_2, right_1, right_2} - {biggest_1})
    return biggest_1, biggest_2

print(biggest_two([10, 20, 30, 35, 2, 17]))
\end{lstlisting}
\begin{verbatim}
(35, 30)
\end{verbatim}
\end{answer}

\begin{subquestion}
find the \href{20221130014441-time_complexity.org}{time complexity} of the algorithm\\
\begin{answer}
we consider the \href{20221105001640-recursive_function.org}{recurrence relation} \(T(n)=2T\left(\frac{n}{2}\right) + C\):\\
\begin{align*}
  T(n) &= 2T\left(\frac{n}{2}\right) + C\\
  T\left(\frac{n}{2}\right) &= 2T\left(\frac{n}{4}\right) + C\\
  T\left(\frac{n}{4}\right) &= 2T\left(\frac{n}{8}\right) + C\\
  T\left(\frac{n}{2^{\lg n}}\right) &= 2T\left(\frac{n}{2^{\lg n+1}}\right) + C\\
  T(n) &= 2\left(2T\left(\frac{n}{4}\right) + C\right) + C = 4T\left(\frac{n}{4}\right) + 3C\\
  &= 4\left(2T\left(\frac{n}{8}\right)+C\right)+3C = 8T\left(\frac{n}{8}\right)+7C\\
  &= 8\left(2T\left(\frac{n}{16}\right)+C\right)+7C = 16T\left(\frac{n}{16}\right)+15C\\
  &\vdots\\
  &= 2^{\lg n}T\left(\frac{n}{2^{\lg n}}\right) + \sum_{i=0}^{\lg n} 2^{i}C\\
  &= 2^{\lg n}C_1 + \sum_{i=0}^{\lg n} 2^{i}C\\
  &= 2^{\lg n}C_1 + \frac{1-2^{\lg(n)+1}}{-1}C\\
  &= 2^{\lg n}C_1-C+2^{\lg(n)+1}\\
  &\Rightarrow \Theta\left(2^{\lg n}\right) = \Theta(n)
\end{align*}
note that although it doesnt matter, \(C_1\) differs a from \(C\) as \(C_1=T(1)\) is the constant time it takes for the base case calls of the recursion\\
\end{answer}
\end{subquestion}
\end{question}

\begin{question}
find the time complexity of the following algorithm\\
\begin{center}
\includesvg{/home/mahmooz/brain/out/h7AJcrd}
\end{center}

\begin{answer}
at first glance this seems like a \(\Theta(n^2)\) algorithm but its time complexity is actually \(\Theta(n)\)\\
we track the value of the variable \texttt{i}: on each iteration in the main loop, the variable \texttt{i} gets incremented as many times as the inner loop runs plus one (\texttt{j} is assigned back to \texttt{i} so we're basically acting on \texttt{i} itself), \texttt{i} can be at most incremented \texttt{N} times before both loops wouldnt execute anymore, and since on each iteration of either loop \(i\) gets incremented both loops are bound by \texttt{N} iterations at most so the total time complexity is \(\Theta(N)\)\\
\end{answer}
\end{question}

\begin{question}
find the time complexity of the following algorithm\\
\begin{center}
\includesvg{/home/mahmooz/brain/out/yx1C6jB}
\end{center}

on the first iteration the inner loop runs and sets all the elements of the array to the value \texttt{A[i]} and so the \texttt{while} loop would never run again on next iterations of the \texttt{for} loop, so for \texttt{N-1} iterations of the \texttt{for} loop we're doing constant work (all except the first iteration) which amounts to a time complexity of \(\Theta(N)\)\\
adding the \(\Theta(N)\) from the first iteration we get \(\Theta(N) + \Theta(N) \Rightarrow \Theta(N)\)\\
so the time complexity is \(\Theta(N)\)\\
\end{question}

\begin{question}
given a function \texttt{P(A,l,k)} that sorts \texttt{k} elements in \texttt{A} starting at index \texttt{l}\\
the algorithm runs as follows:\\
\begin{enumerate}
\item if k<10, apply \href{20221105005344-sorting_algorithm.org}{bubble sort}\\
\item run \texttt{P(A, l, 3k/4)}\\
\item run \texttt{P(A, l+k/4, 3k/4)}\\
\item run \texttt{P(A, l, 3k/4)}\\
\end{enumerate}
\begin{subquestion}
prove/disprove whether the algorithm actually sorts\\
\begin{answer}
we prove the algorithm using \href{20220707193301-mathematical_induction.org}{perfect induction}:\\
for all \(k < 10\), the algorithm indeed works as we know bubble sort works\\
we assume that the algorithm works for all \(n > k \ge 10\)\\
we assume descending order of sorting \textbf{without loss of generality}\\
we prove the algorithm works for \(k=n\):\\
when \(n=k\) this is what the algorithm does:\\
\begin{enumerate}
\item \(A\left[l,\dots,l+\frac{3n}{4}-1\right]\) is sorted\\
\item \(A\left[l+\frac{n}{4},\dots,l+n-1\right]\) is sorted\\
before this second sorting operation, we knew for sure that the elements in the first fourth of the array are all smaller than the elements in the following 2/4's of the array because this 3/4 subarray as a whole is sorted\\
after the sorting operation the last fourth of the array would be merged into the 2/4's in the middle resulting in the last 3/4's in the array being sorted, now we knew before that 2/4's of the elements in the array are bigger than the elements in the first fourth, so after merging the last fourth we now know that there are \textbf{atleast} 2/4's that are bigger than the first fourth:\\
\[ \left(\forall x \in A\left[l,\dots,l+\frac{n}{4}-1\right]\right)\left(\forall y \in A\left[l+\frac{2n}{4},\dots,l+n-1\right]\right)[x \ge y] \]\\
in words: each element in the last 2/4's of the array is bigger than all elements in the first fourth of the array\\
\item sort \(A\left[l,\dots,l+\frac{3k}{4}-1\right]\)\\
a consequence of the third sorting step is that the first fourth would be merged into the the second 3/4's and its elements would be spread across the first 2/4's, we know the last 2/4's wouldnt be affected by this sorting operation because they are all for sure bigger than the first fourth\\
and so the array would be for sure sorted after the third step.\\
\end{enumerate}
and since we've proved that the algorithm is correct for \(k=n\) we've proved its correctness by perfect induction\\
\end{answer}
\end{subquestion}
\begin{subquestion}
find the time complexity of the algorithm using \href{20221105001640-recursive_function.org}{recursion tree}s, find the number of nodes in each level and the number of levels\\
\begin{note}
im almost sure there is something wrong here but im not sure how to improve this solution\\
\end{note}
\begin{answer}
we propose the following \href{20221105001640-recursive_function.org}{recurrence relation}:\\
\[
  T(n) = 3T\left(\frac{3}{4}n\right) + C
\]\\
and build a recursion tree from it:\\
\begin{center}
\includesvg{/home/mahmooz/brain/out/hzown4g}
\end{center}

let \(x\) denote the height of the tree, we know \(\left(\frac{3}{4}\right)^xn < 10\) because \(n < 10\) is the \href{20221105001640-recursive_function.org}{base case}\\
the number of nodes in each level \(y\) is \(3^{y-1}\)\\
according to the pattern in the tree:\\
\[
T(n) = \sum_{y=0}^{x-1} 3^yC
\]\\
we might have to check the time complexity for each value in \(T(1),\dots,T(9)\) assigned as the base case only if they do play a role in the total time complexity\\
the time complexity is the sum of all the nodes in the tree:\\
\begin{gather*}
  \left(\frac{3}{4}\right)^xn = 9\\
  \ln\left(\left(\frac{3}{4}\right)^xn\right) = \ln9\\
  \ln{n} + \ln\left(\frac{3}{4}\right)^x = \ln9\\
  x = \frac{\ln9 - \ln n}{\ln\left(\frac{3}{4}\right)}\\
  T(n) = \sum_{y=0}^{x-1} 3^yC = \left(\frac{3^x}{2} - \frac{1}{2}\right)C = \left(\frac{3^\frac{\ln9 - \ln n}{\ln\left(\frac{3}{4}\right)}}{2} - \frac{1}{2}\right)C\\
  \Theta\left(\left(\frac{3^\frac{\ln9 - \ln n}{\ln\left(\frac{3}{4}\right)}}{2} - \frac{1}{2}\right)C\right) \implies \Theta\left(3^{\ln n}\right) \implies \Theta(3^n)
\end{gather*}
\begin{note}
it confused me for a good while that we have \(-\ln n\) and i thought that the expression would converge as \(n\) goes to infinity but then i realized that \(\ln\frac34\) is negative and it would cancel the minus sign so the time would actually diverge\\
\end{note}
we can notice that the constant \(9\) played no part in the equation so the result would be the same for all constants 1 to 9\\
so the time complexity is \(\Theta\left(3^n\right)\)\\
the height of the tree is dependant on the initial value of \(n\) but i approximate it to be \(\frac{\ln9 - \ln n}{\ln\left(\frac{3}{4}\right)}\)\\
\end{answer}
\end{subquestion}
\end{question}

\begin{question}
modify the \texttt{merge} algorithm such that it merges without using elements that have a value of \(\infty\)\\
\begin{answer}
\begin{center}
\includesvg{/home/mahmooz/brain/out/dZUUjjV}
\end{center}
\end{answer}
\end{question}
\end{document}